\documentclass{article}
\usepackage[utf8]{inputenc}
\usepackage[T1]{fontenc}
\usepackage[a4paper]{geometry}		% Reduire les marges
\geometry{hmargin=3cm,vmargin=3.5cm}
\usepackage[english]{babel} 
\usepackage{amsfonts}
\usepackage{amsthm}
\usepackage{amsmath}
\usepackage{amssymb}
\usepackage{hyperref}
\usepackage{array}
\usepackage[svgnames]{xcolor}
\usepackage{xpatch}
\usepackage{tikz}
\usepackage{mathrsfs}
%\usepackage{graphicx}
%\usepackage{caption}
\usepackage[svgnames]{xcolor}
\usepackage{eurosym}
\usetikzlibrary{positioning,calc}
\renewcommand{\thesection}{\Roman{section}}
\xpatchcmd{\proof}{\itshape}{\normalfont\proofnameformat}{}{}
\newcommand{\proofnameformat}{\bfseries}
\theoremstyle{definition}
\newtheorem{defi}{Definition}[section]
\newtheorem{theo}{Theorem}[section]
\newtheorem{prop}[theo]{Proposition}
\newtheorem{lem}[theo]{Lemma}
\newtheorem{coro}[theo]{Corollary}
\newtheorem{exe}{Example}[section]
\newtheorem{rem}{Remark}[section]
\newtheorem{nota}{Notation}[section]
\renewcommand\qedsymbol{$\blacksquare$}
\numberwithin{equation}{section}

\DeclareMathOperator{\R}{\mathbb{R}}
\DeclareMathOperator{\Z}{\mathbb{Z}}
\DeclareMathOperator{\Q}{\mathbb{Q}}
\DeclareMathOperator{\N}{\mathbb{N}}
\DeclareMathOperator{\E}{\mathbb{E}}
\DeclareMathOperator{\C}{\mathcal{C}}

\title{Exo 5, correction}       
\author{Guillaume Woessner}
\date{}                       % sans date : celle du jour de compilation

%\pagestyle{headings}          % Pour mettre des entetes avec les titres
                              % des sections en haut de page
\sloppy                       % Ne pas faire deborder les lignes dans la marge

\begin{document}

\maketitle                    % Faire un titre utilisant les donnees
                              % passees : \title, \author et \date

\begin{abstract}
\begin{center}
Un peu de tout.
\end{center}
\end{abstract}

%\tableofcontents              % Table des matieres


Dans tout ce document, l'espace de probabilité considéré est $(\Omega,\mathcal{F},(\mathcal{F}_n)_n,\mathbb{P})$.



\paragraph{Exercice 1} [Révisions]~\\
Soit $(S_n)_{n\in \N}$ une marche aléatoire sur $\Z$, $\tau:=\tau_{-2,1}=\min\{n\in\N~:~S_n=1 \text{ ou } S_n=-2\}$, et $\tau':=\tau_1 = \min\{n\in\N~:~S_n=1\}$. Rappelez brièvement pourquoi $\tau<\infty$ presque sûrement et $\E[\tau]=2$, et en déduire que $\tau'<\infty$ presque sûrement, et que $\E\tau'=\infty$.

\begin{proof}
Pour montrer que $\tau$ est presque sûrement fini, on définit les évènements $A_n:=\{X_i=1~\forall 3n+1 \leq i \leq 3(n+1)\}$ et en remarquant que si l'un au moins des $A_n$ se réalise alors $\tau$ est fini, on obtient que $\tau \leq 3\tilde{\tau}$ où $\tilde{\tau}$ est le rang du premier $A_n$ qui se réalise. Or celui-ci est presque sûrement fini par Borel-Cantelli.\\
Pour montrer que $\E\tau<\infty$ on applique le théorème d'arrêt au temps d'arrêt borné $\tau\wedge n$ et on fait tendre $n$ vers l'infini.\\
~\\
On obtient alors directement que $\E\tau'=\infty$ en remarquant que pour tout $M>0$, on a $\tau'\geq \tau_{-M,1}$ et $\E\tau_{-M,1}=M$.\\
De même on sait que $S_{\tau_{-M,1}}=1$ avec probabilité $\frac{M}{M+1}$ et $S_{\tau_{-M,1}}=-M$ avec probabilité $\frac{1}{M+1}$. On a donc $\mathbb{P}(\tau'=\infty)\leq \mathbb{P}(\tau_{-M,1}=\infty)+\mathbb{P}(S_{\tau_{-M,1}}=-M) = 0 + \frac{1}{M+1}\mapsto 0$.
\end{proof}


\paragraph{Exercice 2} 
Soit $(S_n)_{n\in \N}$ une sur-martingale, et $\tau$ un temps d'arrêt. Montrez que le processus $(S_{\tau\wedge n})_{n\in \N}$ est encore une sur-martingale.

\begin{proof}
Il existe deux façons de résoudre cet exercice. Soit on peut reprendre la preuve du Lemme 3.5 affirmant qu'une martingale arrêtée reste une martingale. Soit on peut utiliser la décomposition de Doob, et c'est celle-ci qu'on va exposer ici.\\
D'après le théorème de décomposition de Doob, on sait qu'il existe des uniques processus $D_n$ croissant prévisible et $M_n$ martingale tels que $S_n = M_n + D_n$. Ainsi $S_{\tau\wedge n} = M_{\tau\wedge n} + D_{\tau\wedge n}$. Or d'après le Lemme 3.5 $M_{\tau\wedge n}$ reste une martingale, et bien sûr $D_{\tau\wedge n}$ reste croissant prévisible. Donc $S_{\tau\wedge n}$ reste bien une sur-martingale.
\end{proof}


\paragraph{Exercice 3}[Formule de Wald]~\\
Soit $(X_n)_{n\in\N^\times}$ une suite \textit{iid} de variables aléatoires intégrables adaptée à la filtration et $\tau$ un temps d'arrêt intégrable adapté à la filtration. Supposons en outre que la tribu $\mathcal{F}_n$ et la variable $X_{n+1}$ sont indépendantes. Explicitez la définition de $S:=\sum_{n=1}^\tau X_n$, et montrez que
$$\E[S]=\E[X_1]\E[\tau].$$

\begin{proof}
$S$ est la variable aléatoire $\sum_{n=1}^\infty X_n \mathbf{1}_{1\leq n\leq\tau}$.\\
De plus $\{1\leq n\leq\tau\}=\{\tau\leq n-1\}^c \in \mathcal{F}_{n-1}$, et donc $X_n$ est indépendante de $\mathbf{1}_{1\leq n\leq\tau}$. Ainsi $S$ est intégrable car 
$$\E|S|\leq \sum_n \E|X_n|\E\mathbf{1}_{1\leq i\leq\tau} = \sum_n \E|X_1|\mathbb{P}(1\leq i\leq\tau)=\E|X_1|\sum_n\mathbb{P}(1\leq i\leq\tau)=\E|X_1|\E[\tau].$$
On est alors justifiés à effectuer les mêmes calculs pour démontrer la formule de Wald.
\end{proof}


\paragraph{Exercice 4}
Trouver une martingale $(X_n)_{n\in \N}$ et un temps d'arrêt $\tau$ tels que $X_\tau$ soit non intégrable.

\begin{proof}
Soit $(X_n)_{n\in \N}$ une marche aléatoire sur $\Z$, soit $Y$ une variable aléatoire indépendante de $(X_n)_{n\in \N}$, et soit $\tau:=\inf\{n\in\N~:~X_n = Y\}$. La variable aléatoire $X_\tau$ convient. % Y should be non-integrable
\end{proof}


\paragraph{Exercice 5}
Trouver deux martingales $(X_n)_{n\in\N^\times}$ et $(Y_n)_{n\in\N^\times}$, ainsi qu'un temps d'arrêt $\tau$ par rapport à la filtration tels que le théorème d'arrêt s'applique à $X_\tau$ mais pas à $Y_\tau$.

\begin{proof}
Soient $(X^0_n)_{n\in\N^\times}$ et $(Y_n)_{n\in\N^\times}$ deux marches aléatoires sur $\Z$ indépendantes, et $\tau:=\inf\{n\in\N~:~Y_n = 1\}$. Soit ensuite $X_n:=X^0_{n\wedge \nu}$ avec $\nu:=\inf\{n\in\N~:~|X^0_n|=10\}$, \textit{ie} $(X_n)_{n\in\N^\times}$ est la marche aléatoire $(X^0_n)_{n\in\N^\times}$ arrêtée en $\pm 10$. On peut vérifier que $(X_n)_{n\in\N^\times}$, $(Y_n)_{n\in\N^\times}$ et $\tau$ définis de cette manière répondent à la question.\\
\textbf{PS} : on aurait pu même aller encore plus vite en posant plutôt $(X_n)_{n\in\N^\times}\equiv 0$.
\end{proof}


\paragraph{Exercice 6}
Calculer le temps moyen d'attente pour lire votre date de naissance dans une suite de nombre générée aléatoirement, à raison d'un nouveau nombre par minute.

\begin{proof}
On raisonne avec la date 01012001, mais le raisonnement pourra un peu changer en fonction de votre date de naissance.\\
Avant tout il faut définir une martingale. On considère la situation suivante, dans lequel vous êtes propriétaire d'un casino pour lequel à chaque minute un nouvel individu vient parier avec vous de la manière suivante,
\begin{itemize}
\item[$\bullet$] Il paie 1 franc pour jouer, et parie sur le fait que le premier chiffre qui sortira est un 1. S'il gagne il remporte 10 francs, sinon vous gagnez sa mise.
\item[$\bullet$] Si vous avez gagné, il repart bredouille. Sinon il remise les 10 francs gagnés sur le 2. S'il gagne il empoche 100 francs, sinon vous gagnez sa mise de départ.
\item[$\dots$]
\item[$\bullet$] Tant qu'il gagne, il continue et remise l'entièreté de ses gains sur le prochain chiffre, que vous multipliez par 10 s'il gagne, sinon il perd tout.
\item[$\bullet$] S'il gagne au dernier coup (le huitième) et qu'ainsi il a réussi à prédire l'occurence de votre date de naissance, il repart avec $10^8$ francs. Félicitations à lui, cet évènement avait une probabilité $10^{-8}$ d'occurer Dans le cas contraire vous gagnez son franc de départ.
\end{itemize}
On peut maintenant spécifier la martingale considérée. Soit $S_n$ vos profits à la minute $n$. Par construction $(S_n)_{n\in\N}$ est une martingale par rapport $ (X_n)_{n\in\N}$ la suite des nombres sortant du générateur. Evidemment le temps d'arrêt $\tau$ à considérer est le premier temps pour lequel quelqu'un gagne à ce jeu. Finalement on introduit également les notations $S_n = G_n - P_n$ avec $G_n$ vos gains au temps $n$ et $P_n$ vos pertes. Notez bien que par définition $G_n=n$, car chaque joueur paie 1 franc pour jouer.\\
~\\
Premièrement, il nous faut montrer que le théorème d'arrêt s'applique, \textit{ie} que $\E[S_\tau]=\E[S_0]=0$. Considérons d'abord le temps d'arrêt $\tau\wedge n$ borné pour un entier naturel $n$. On peut ainsi appliquer le théorème d'arrêt à ce nouveau temps d'arrêt et obtenir $\E[S_{\tau\wedge n}]=\E[S_0]=0$. Rappelons que $\E[S_{\tau\wedge n}]=\E[G_{\tau\wedge n}]-\E[P_{\tau\wedge n}]$ et notez que $\E[G_{\tau\wedge n}] = \E[\tau\wedge n] \mapsto_n \E[\tau] = \E[G_\tau]$ par le théorème de convergence monotone. De même $\E[P_{\tau\wedge n}]\mapsto_n \E[P_\tau]$ par le théorème de convergence dominée ($P_{\tau\wedge n}$ est dominée par $8\times 10^8$. Ainsi $\E[S_{\tau\wedge n}]\mapsto_n\E[S_\tau]$.\\
Finalement on a obtenu $\E[\tau]=\E[G_\tau]=\E[P_\tau]$. Or, au temps $\tau$, les deux seuls individus à avoir gagné quelque chose sont celui qui a gagné au jeu (dont les gains sont de $10^8$ francs), et celui qui est arrivé juste avant le dernier 0 (dont les gains sont de $10^2$ francs). Ainsi $\E[\tau]=\E[L_\tau]=10^2+10^8$.
\end{proof}











\end{document}
